
\begin{frame}{Análise da Distribuição dos Pesos das Arestas}
    \begin{block}{Distribuição dos Pesos}
        \begin{itemize}
            \item \textbf{Distribuição dos Pesos das Arestas:} 
            \begin{itemize}
                \item Foco na análise de como os pesos das arestas (hierarquia das \textit{medalhas}) estão distribuídos na rede de atletas.
            \end{itemize}
            \item \textbf{Função de Distribuição Cumulativa (CDF):}
            \begin{itemize}
                \item Mostra a \textit{probabilidade acumulada} de um peso de aresta ser \textbf{menor ou igual} a um valor.
                \item Eixo X: Peso da Aresta; Eixo Y: Probabilidade Acumulada.
                \item Curva ascendente: indica a proporção de arestas com pesos menores que um dado valor.
            \end{itemize}
        \end{itemize}
    \end{block}
\end{frame}

\input{contents/AnaliseDistribuiçãoComulativa FUTEBOL}

\input{contents/AnaliseDistribuiçãoComulativa BQTBL}

\input{contents/AnaliseDistribuiçãoComulativa NTÇ}

\begin{frame}{Resumo da Análise das CDFs e Validação da Modelagem}
    \begin{block}{Validação da Modelagem de Pódio Esportivo}
        \begin{itemize}
            \item As CDFs revelam que a maioria das arestas possui pesos baixos, refletindo pequenas diferenças de desempenho entre os atletas, o que é consistente com a proximidade competitiva esperada.
        \end{itemize}
    \end{block}
    
    
\end{frame}

